\section{Simulation-Based Inference for $dN/dS$}
\label{sec:sbi_cnn}

The 1pPDF formalism reviewed in the previous section represents the most developed analytical approach to sub-threshold source counts.
It has produced the strongest existing measurements of the gamma-ray $dN/dS$ below the \textit{Fermi}-LAT detection limit.
However, its underlying assumptions --- a parametric model for the source-count distribution, a simplified treatment of spatial correlations, and a computationally expensive likelihood --- impose limitations that motivate a complementary approach.

This section discusses those limitations in detail (Section~\ref{sec:6.3.1}), then introduces the simulation-based inference framework that addresses them (Section~\ref{sec:6.3.2}), and finally summarizes the handling of the spherical geometry of the input data (Section~\ref{sec:6.3.3}).

\subsection{Limitations of the Analytical Framework}
\label{sec:6.3.1}

The generating function formalism underlying the 1pPDF method rests on several simplifying assumptions, each of which constrains its applicability.

\textit{Neglect of spatial correlations.}
The 1pPDF treats each pixel independently.
In practice, however, the point-spread function (PSF) of the \textit{Fermi} LAT smears a point source across multiple pixels, introducing spatial correlations that the generating function approach does not capture.
These correlations carry information about the source-count distribution that is discarded in the single-pixel treatment, reducing the effective constraining power and potentially introducing biases.

\textit{Parametric dependence.}
The $dN/dS$ is typically modeled as a multiply broken power law with a fixed number of segments.
This assumption is sufficiently flexible for smooth, featureless distributions, but it cannot accommodate spectral features, bumps, or breaks at unforeseen flux values.
The reconstruction is therefore conditional on the assumed functional form.

\textit{Computational cost and absence of amortization.}
The explicit 1pPDF likelihood is expensive to evaluate: each call requires recomputing the generating function at every pixel.
Sampling the posterior with MCMC compounds this cost, since exploring the high-dimensional space of $dN/dS$ and nuisance parameters (such as the Galactic foreground normalization) requires many such evaluations.
Moreover, the inference is not amortized: every new sky map, whether observed or simulated, demands a fresh set of MCMC chains.
This becomes a concrete disadvantage when the method must be validated, since quantifying potential biases requires running the analysis on simulated maps with known $dN/dS$, each carrying the full cost of new chains.
In practice, achieving a robust validation therefore comes at a substantial computational price.

\textit{Limited multi-band extension.} While multi-energy extensions of the 1pPDF have been performed \cite{Zechlin:2016pme,Lisanti:2016jub}, extending the formalism to a large number of energy bins simultaneously is practically challenging, since the number of model parameters and the computational cost grow rapidly.

The community has recognized some of these issues: Zechlin et al.
\ developed a Compound Poisson Generator (CPG) formalism \cite{Cuoco-1pdf} to address potential biases in the 1pPDF approach, though this alternative still operates within the generating-function framework.
A fundamentally different strategy is needed to overcome all four limitations simultaneously.

\subsection{A Simulation-Based Approach}
\label{sec:6.3.2}

The limitations identified above share a common root: the 1pPDF relies on an a slow analytic likelihood that incorporates only a subset of the information contained in the photon-count map.
An alternative is to bypass the likelihood entirely and employ implicit likelihood methods, the general framework introduced in Section~\ref{sec:3.2}.

\blue{As discussed in Chapter~\ref{ch:methods}, in the \SBI paradigm the likelihood is never evaluated explicitly: it is defined implicitly by the simulator (Section~\ref{sec:3.2.1}).
In general, writing the likelihood would require marginalizing over all the latent variables of the simulator; sampling from it, by contrast, only requires running it forward.
In practice,} one draws source fluxes from a parametric $dN/dS$, places them on the sky following an assumed spatial distribution, convolves with the instrument PSF and exposure, adds the Galactic foreground and isotropic background, and applies Poisson noise to produce a synthetic photon-count map.
Each simulation is fast (around 0.1s) and faithful to the known instrumental and astrophysical effects.
\blue{Crucially, the latent variables --- here, the position and flux of every individual source --- never need to be handled explicitly: each synthetic map is one random realization of them, so the intractable marginalization is performed implicitly by the simulator itself.}

A convolutional neural network (CNN) is then trained on an ensemble of such simulated maps to learn the mapping from a photon-count sky to the underlying $dN/dS$, using the architecture and training methodology detailed in Section~\ref{sec:3.3}.
% The specific configuration adopted for the $dN/dS$ problem discretizes the output into 20 flux bins spanning $[5 \times 10^{-12},\, 10^{-7}]$~cm$^{-2}$~s$^{-1}$, plus the isotropic background level $F_\mathrm{iso}$, yielding 21 output parameters.
% An ensemble of $9 \times 10^5$ synthetic maps is generated for training.
\blue{The configuration adopted for the $dN/dS$ problem discretizes the output into a set of flux bins plus the isotropic background level $F_\mathrm{iso}$; the exact binning, flux range, and training-ensemble size are specified later in this chapter (Section~\ref{sec:architechture-and-training}).}
% The discretized representation is critical: unlike the parametric power-law adopted by the 1pPDF, it allows the CNN to reconstruct the $dN/dS$ non-parametrically, without assuming a fixed functional form.
The network employs the heteroscedastic Gaussian cost function and concrete dropout regularisation to estimate both the $dN/dS$ and its uncertainties; the full formalism is presented \blue{later in this chapter} (Section~\ref{sec:bayesian-error}). A frequentist cross-check using the distribution of residuals on held-out validation maps confirms that the Bayesian error bands are well calibrated (Section~\ref{sec:bias}).

This framework addresses the four limitations identified in Section~\ref{sec:6.3.1}: inter-pixel PSF correlations are captured implicitly by the convolutional architecture, the $dN/dS$ is recovered non-parametrically, inference is amortized (additionally a trained network produces predictions in milliseconds), and the pipeline extends naturally to multiple energy bins by adding input channels.

\subsection{Inference on Spherical Data}
\label{sec:6.3.3}

\textit{Fermi}-LAT data are represented as HEALPix sky maps on the sphere.
Applying standard planar CNNs to such data requires projecting the sphere onto a flat image, which inevitably introduces distortions.
The \texttt{map2patches} algorithm, described in Section~\ref{sec:3.3}, addresses this by subdividing the HEALPix sphere into 12 equal-area base patches at order $n=0$, each of which is independently mapped to a 2D image and processed through standard 3D convolutions.
This approach preserves the area and information content of each pixel without spherical distortion, while leveraging the optimized performance of standard convolutional operations.

A fully spherical cross-check using the NNHealpix architecture \cite{Krachmalnicoff:2019} yields compatible results, confirming that the map2patches approximation is reliable for the isotropic, small-scale signal of interest.
The map2patches implementation is approximately an order of magnitude faster, enabling efficient generation of the large training ensembles required by the SBI framework.
Full architectural details are presented \blue{later in this chapter} (Section~\ref{sec:architechture-and-training}).
